\documentclass[12pt]{article}
\usepackage[margin=0.7in]{geometry}

% \usepackage[margin=0.8in]{geometry}

% \usepackage{newpxtext}
% \usepackage{newpxmath}

\usepackage[tt=false]{libertine}

% \usepackage{newtxtext}
% \usepackage{newtxmath}

\usepackage[stretch=10]{microtype}

\linespread{1.05}

\usepackage{enumitem}
\setlist{nosep}

\usepackage[hidelinks, linktoc=all]{hyperref}

% Answers to the review questions. In the PDF this typesets the answer as an
% indented block right after the question; on the web the same environment is
% turned into a click-to-reveal dropdown (see web-pipeline/qa-details.lua).
% Pandoc does not implement \NewDocumentEnvironment, so it leaves the environment
% as a Div the filter can restructure, while LaTeX honours the definition here.
\NewDocumentEnvironment{answer}{}
  {\par\smallskip\noindent\ignorespaces}
  {\par\smallskip}

\title{Ethics}
\date{}
\author{}


\begin{document}

\section{Metaethics}
The study of the foundation of morality.

\noindent Theme: Debate over whether moral facts exist universally and independently of human minds.
Is there a single standard of Right and Wrong that applies to human beings everywhere?

\subsection{Cognitivism vs Non-cognitivism}
Are moral statements like "stealing is wrong" stating a fact about the world,
or just expressing one's personal feelings?

\textbf{Cognitivism} This view believes that moral statements are truth apt,
i.e. they can be objectively True or False. Cognitivists have to answer "\textit{what exactly makes these statements true?}".

\textbf{Non-cognitivism} This view argues that moral statements are \textbf{not} capable of 
being True or False. Instead, they are expressions of our inner emotions, attitudes, or personal preferences.
Note that being a non-cognitivist does not make one a moral skeptic (rejecting morality altogether).

\subsection{Emotivism}
Emotivism is a specific type of non-cognitivism, which claims that moral statements are expressions of emotions or feelings.
"\textit{Stealing is wrong}" is expressing our emotion of disapproval.

\subsection{Error Theory \& Extreme Nihilism}

\textbf{Error theory} is a form of cognitivism. It agrees moral statements \textit{attempt} to express objective facts.
But error theorists claim that all ethical convictions are false beliefs.
For the statement "\textit{Stealing is wrong}", Error Theory argues this is making a claim about a property ("wrongness")
that does not exist in the world. Hence, any moral statement is technically a false belief.

\textbf{Extreme Nihilism} extends Error theory to everyday lives. It claims that nothing actually has value, i.e. the concepts of right vs wrong,
good vs bad, simply do not exist. The entire concept of morality is an illusion.
The argument:
\begin{enumerate}
    \item We should only believe in something (or accept we have evidence for it) if that thing helps us explain the physical world we observe around us.
    For example, we believe in gravity because it explains why an apple falls to the ground.
    \item "Moral hypotheses" (like the idea that an action was "evil" or "unjust") never actually help us explain why anything happens in the physical world. 
    Because moral facts don't explain physical observations, we have absolutely no evidence that our moral opinions are true.
\end{enumerate}
This means that there are no moral constraints in the world (right and wrong are illusions), so we can do anything we want (no rules governing our behaviour).

\subsection{Moral Realism}
Moral realism is a form of cognitivism.

\noindent \textbf{Claim}: There are objective, mind-independent moral facts.
Mind-independent means that the moral facts remain objectively and completely true, similar to math or physical facts,
regardless of what some individuals or societies think or feel.
The statement "\textit{Stealing is wrong}" is a factual, universal truth regardless of what some individuals may think.

\noindent Why might one be a moral realist:
\begin{enumerate}
    \item \textbf{Establish secular objectivity} Provides a basis for moral truths that exist regardless of cultural or personal beliefs.
    For example, "\textit{Murder is wrong}" is a natural fact of the universe that we can discover, just like gravity or biology.
    \item \textbf{Validate Human Intuition} Moral realism allows us to do justice to our innate sense of right and wrong.
    It validates our strong intuition that ethical issues are genuine facts, rather than mere matters of opinion.
    For example, "\textit{Torturing an innocent puppy for fun is wrong}", moral realism would validate our feelings that that act is factually, objectively wrong.
\end{enumerate}

\subsection{Ethical Naturalism}
A version of moral realism (right and wrong are objective facts).
Claims:
\begin{enumerate}
    \item \textbf{Moral realism} There are objective, mind-independent moral facts.
    \item \textbf{Metaphysical Naturalism} The belief that all facts are natural facts.
    (Since all substances and properties that exist are purely natural, then nothing supernatural exists).
    Moral properties are strictly defined by facts about the natural, physical world.
    The "wrongness" of murder is just a complex but natural fact, perhaps rooted in biological survival etc.
    \item \textbf{Epistemic Naturalism} Claims that because moral facts are just natural facts,
    we can discover and verify them using the same methods in the natural sciences like physics.
    For example, to know if something is "morally good", we may observe its physical and psychological effects on humans.
\end{enumerate}
Objective right and wrong exists, and are completely natural, physical facts, therefore we can use science, observation, logic to figure out what those moral facts are.

\vspace{1em}
\noindent\textbf{Example of Ethical Naturalism} Hedonistic Utilitarianism: It claims that 'Goodness' is exactly equal to 'Pleasure.'
Because pleasure is a biological, psychological, and entirely natural phenomenon that can be scientifically observed, this theory successfully grounds a moral concept in the physical, natural world.

\subsection{Ethical Supernaturalism}
Is a form of moral realism, agreeing that there are objective, mind-independent moral facts.
However, ethical supernaturalists claim that ethical propositions are grounded in a divine or supernatural realm, rather than in nature.

\textbf{Divine Command Theory} Some supreme being, perhaps God, is the some ultimate source of morality. An action is morally good strictly because it is commanded by a divine being.


\subsection{Two Objections Against Ethical Naturalism}
\subsubsection{Science vs Ethics (Harman's Challenge)}
Recall that ethical naturalism claims that we can use scientific methods to discover moral facts.
Gilbert Harman's counter-argument asserted that moral facts are unnecessary to explain our observations.

\noindent A scientific example: When a physicist looks into a cloud chamber and observes a proton vapour trail,
the existence of the physical proton is strictly necessary to explain that observation.

\noindent A moral example: When we see someone stealing a homeless person's wallet,
we observe that "\textit{that's wrong}". The only reason for this observation is our own psychology, upbringing etc.
Note that we did not anything physical, objective property called "wrongness" to explain that observation.

\noindent Science relies on Ockham's Razor: we should only believe in things that are strictly necessary to explain our observations.
Objective, moral facts are not necessary to explain our moral observations.
Since ethical naturalism uses scientific standards to claim the existence of moral facts,
the same scientific standards disproved the existence of moral facts.

\textbf{Counter to Harman} Mathematical proofs and equations are widely accepted
as objectively true, although they do not require empirical, physical observation. 
E.g. in equations like 5 + 7 = 12, We can't physically perceive or touch the number 12.
In the same way, it is impossible to physically perceive abstract moral qualities.
Therefore, we should accept moral facts just as we accept something like mathematics,
as not all objective facts have to be physical things we can observe. 

\subsubsection{Moore’s Open Question Argument}
Ethical naturalism claims that moral concepts like "Goodness" are just natural properties like "pleasure" or "biological survival".
G.E. Moore argues that identifying a moral property with a natural property is a fallacy.

\noindent A closed question is a tautology, like "\textit{Is a bachelor an unmarried man?}".

\noindent An open question implies that the answer is not contained within the definitions of the terms alone.
For any proposed natural definition of "good", (like "pleasure"), the question
"\textit{Is pleasure good?}" would be a tautology under ethical naturalism.
However, asking if something is good is always an open, meaningful question.
Therefore, Moore argues that "Goodness" is an unanalysable, non-natural property that cannot be defined by anything else. 

\noindent\textbf{Counter to Moore} Meaning is not the same as synthetic identity.
Moore's argument relies on analytic identity, that for two things to be identical,
their words must have the exact same dictionary meaning.
However, there is also synthetic identity, like "\textit{Water is H20}", 
where it might have been an open question hundreds of years ago, but after scientic discovery, the two concepts are actually the same thing in the real world,
even if their dictionary definitions are different.

\noindent Hence, questions like "\textit{Is pleasure good?}" may be an open question today,
but it may mean that we are still in the process of scientifically discovering the synthetic identity of "Goodness". Moral properties could be natural properties even if their concepts are different.

\subsection{Reductive Naturalism}
Reductive naturalism presents a solution to Harman's challenge that we cannot physically observe moral facts.
Reductive naturalism argues that we can construct moral facts out of other natural, observable facts.
Scientific example: Colour can be broken down into a combination of facts about light waves, physical surfaces and human eyesight.
Similarly, "goodness" is simply a combination of biological, psychological and social facts.
Reductive naturalists use Functionalism to do this, which is the idea that something is good if it effectively fulfills its purpose.
For example, a watch is "good" if we can observe that it keeps accurate time, if it's durable, if it fits well, etc.
We can apply this to human beings, where a "good" human action might be one that adequately fulfills the function of promoting human survival,
societal harmony, or psychological well-being. 

\noindent\textbf{The problem with Functionalism} is that humans do not have a single, objective, universally agreed-upon "purpose" or "function".
For example, if a naturalist claims a human's purely natural function is "biological reproduction and survival," does that mean a person who chooses to be childless or sacrifices their life for a stranger is functionally "broken" or "morally bad"?
Reducing human morality to biological functions often leads to highly controversial and counterintuitive conclusions.

\end{document}
